% hci-spine LaTeX skeleton — FORM: artifact (system / tool / technique / design artifact)
% Venues: UIST, CHI, DIS, TEI. Tradition modifiers: systems-building (teaser+tech eval),
% research-through-design (annotated portfolio + strong concepts — uncomment that block).
\documentclass[sigconf,review,anonymous]{acmart}
\settopmatter{printacmref=false}
\begin{document}
\title{TITLE — name the artifact/technique (for TEI, signal the body/material)}
\begin{abstract}
% Problem, the artifact, headline eval numbers. Universal check: name ALL contributions.
\end{abstract}
\maketitle

\section{Introduction}            % problem + "WHY IS THIS HARD" + explicit contribution list
\section{Related Work}            % technique LINEAGE, not a topic dump
\section{Design Goals / Rationale}
\section{System / Technique}       % the contribution
\begin{figure}[t]\centering\caption{Teaser: the artifact in one glance.}\end{figure}
\section{Walkthrough / Usage Scenario}
\section{Evaluation}               % technical and/or usability and/or demonstration
% Honesty: every number traces to a log/study; no phantom baseline; real outputs (label mockups/WoZ).
% Media-fit lens (artifact): why not paper/plain screen? AI removable? interaction makes meaning?
%% research-through-design tradition (uncomment if used):
%% \section{Annotated Portfolio / Strong Concepts}  % transferable intermediate knowledge
\section{Discussion}
\section{Limitations}
\section{Conclusion}
\bibliographystyle{ACM-Reference-Format}\bibliography{refs}
\end{document}
